% TOP_PROBLEM: {"id": 171, "title": "Pyjama Set Covering", "category_id": 6, "category": {"id": 6, "name": "geometry", "display_name": "Geometry", "description": "Euclidean and non-Euclidean geometry, geometric structures.", "slug": "geometry", "order_index": 6, "created_at": "2026-07-31T15:26:25.671Z"}, "status": "open"}
% FIRST_POSTED: null
% ENTRY_KIND: statement_only
% Imported from: unsolvedmath_additional_500.tex; UnsolvedMath ID 171
% !TeX program = lualatex
% Compile twice: lualatex 171.tex
% Self-contained: no external bibliography, images, or \input files.
% Numeric section labels and visible numbers are original UnsolvedMath IDs.
\documentclass[11pt,a4paper]{article}
\usepackage[margin=24mm,headheight=14pt]{geometry}
\usepackage{fontspec}
\setmainfont{Latin Modern Roman}
\setsansfont{Latin Modern Sans}
\setmonofont{Latin Modern Mono}
\usepackage{amsmath,amssymb,mathtools}
\usepackage{unicode-math}
\setmathfont{Latin Modern Math}
\usepackage{microtype}
\usepackage{enumitem,xurl,needspace}
\usepackage[dvipsnames]{xcolor}
\usepackage{hyperref}
\hypersetup{unicode=true,colorlinks=true,linkcolor=MidnightBlue,urlcolor=MidnightBlue,
 pdftitle={UnsolvedMath Problem 171},pdfauthor={Source-annotated compilation},bookmarksnumbered=true}
\usepackage{bookmark,tocloft,titlesec,fancyhdr}
\setlength{\cftsecnumwidth}{4.2em}
\cftsetpnumwidth{2.5em}
\cftsetrmarg{4em}
\titleformat{\section}{\Large\bfseries\raggedright}{\thesection}{0.7em}{}
\pagestyle{fancy}\fancyhf{}
\fancyhead[L]{\small\sffamily Additional UnsolvedMath statements}
\fancyhead[R]{\small Original catalogue IDs}
\fancyfoot[C]{\thepage}
\setlength{\parindent}{0pt}
\setlength{\parskip}{5pt plus 1pt}
\setlength{\emergencystretch}{3em}
\setcounter{tocdepth}{1}\setcounter{secnumdepth}{1}
\setlist[itemize]{leftmargin=3em,itemsep=4pt,topsep=4pt,parsep=0pt}
\allowdisplaybreaks
\newcommand{\N}{\mathbb{N}}
\newcommand{\Z}{\mathbb{Z}}
\newcommand{\Q}{\mathbb{Q}}
\newcommand{\R}{\mathbb{R}}
\newcommand{\C}{\mathbb{C}}
\newcommand{\F}{\mathbb{F}}
\newcommand{\A}{\mathbb{A}}
\newcommand{\PP}{\mathbb{P}}
\newcommand{\E}{\mathbb{E}}
\newcommand{\eps}{\varepsilon}
\newcommand{\dd}{\,\mathrm d}
\newcommand{\sref}[2]{\hyperlink{source:#1:#2}{[S#2]}}
\newif\ifshowcatalogue
\showcataloguetrue
\newcommand{\cataloguescope}[1]{\ifshowcatalogue
 \paragraph{Statement recorded in the supplied source.}
 #1\fi}
\begin{document}
% UnsolvedMath ID 171; source catalogue code GREEN-083

\setcounter{section}{170}

\section{The Quantitative Pyjama Covering Problem}\label{171}

{\small\sffamily UnsolvedMath ID 171 · Geometry}

\subsection{Definitions and mathematical statement}

\paragraph{Shared notation.}Unless a section says otherwise, $\N=\{1,2,\ldots\}$,
$\Z$ is the ring of integers, $\Q,\R,\C$ are the rational, real, and complex
fields, $[N]=\{1,\ldots,\lfloor N\rfloor\}$, and $\log$ is the natural
logarithm. For real functions of a parameter tending to infinity,
$f\ll g$ and $f=O(g)$ mean $|f|\leq Cg$ eventually for some fixed $C$;
$f\gg g$ reverses this comparison; $f=o(g)$ means $f/g\to0$;
$f\sim g$ means $f/g\to1$; and $f\asymp g$ means both order bounds.
A subscript on $O$ or $\ll$ specifies allowed constant dependence.
The expression $n^{o(1)}$ in an upper bound means growth smaller than
$n^\epsilon$ for each fixed $\epsilon>0$. Local definitions govern any
exception, finite-field convention, or use of the same symbol for another
object. An asymptotic claim is not automatically an effective algorithm.



For $0<\varepsilon<1/2$, let $P_\varepsilon=\{(x,y):\min_{m\in\Z}|x-m|\leq\varepsilon\}$. A permitted motion is a rotation $R_\theta$ about the origin, with no translation. Let $m(\varepsilon)$ be the minimum size of a finite family of angles for which $\bigcup_j R_{\theta_j}P_\varepsilon=\R^2$, taking $m=\infty$ if none exists. Determine this number or its sharp growth as $\varepsilon\downarrow0$. \sref{171}{1} \sref{171}{2}

\paragraph{Statement recorded in the supplied source.}
How many rotated (about the origin) copies of the "pyjama set" $\{(x, y) \in \mathbb{R}^2 : \operatorname{dist}(x, \mathbb{Z}) \leq \varepsilon\}$ are needed to cover $\mathbb{R}^2$? \sref{171}{1}

\paragraph{Residual task reported by the archive.}

The embedded review (2026-08-17) records: Obtain substantially smaller upper bounds, such as epsilon\textasciicircum{}{-C}, and nontrivial matching lower bounds for the minimum covering number. \sref{171}{1}

\subsection{Short English statement}

How many differently rotated families of narrow periodic strips are needed to cover the whole plane?

\subsection{Sources}

{\small

\begin{itemize}

\item[\hypertarget{source:171:1}{S1}] ULAM.AI, \emph{UnsolvedMath}, supplied \texttt{problems.json}, record 171 (GREEN-083), statement and background. The embedded literature-review date is 2026-08-17; that is the archive’s review date, not a fresh certification. Dataset landing page: \url{https://huggingface.co/datasets/ulamai/UnsolvedMath}.

\item[\hypertarget{source:171:2}{S2}] F. Manners, A solution to the pyjama problem, Invent. Math. 202 (2015), 239-270, \nolinkurl{doi:10.1007/s00222-014-0571-7.} \url{https://arxiv.org/abs/1305.1514}. Bibliographic information imported from the supplied record.

\item[\hypertarget{source:171:3}{S3}] N. Kravitz, J. Leng, Quantitative pyjama, arXiv:2510.17744 (2025). \url{https://arxiv.org/abs/2510.17744}. Bibliographic information imported from the supplied record.

\item[\hypertarget{source:171:4}{S4}] B. Green, 100 Open Problems, Problem 41, version. \url{https://people.maths.ox.ac.uk/greenbj/papers/open-problems.pdf}. The author-hosted document was inspected on 27 September 2026; the archive’s local code is not a problem-number locator in this paper.

\end{itemize}

}

\label{end:171}
\end{document}
